\begin{lemma}\label{lem:nonic77-finite-image-Cartier}
Let $Y/\mathbb Q$ be an odd-level modular curve or diamond quotient,
with good reduction at $2$, a rational cusp $c_0$, and finite etale
cuspidal boundary with cyclotomic fields of definition. Suppose that
$T_2,\delta=\langle2\rangle$ descend to its Jacobian $J$ and satisfy
Eichler--Shimura. Let $a$ be an even integer such that
$(T_2-a)J(\mathbb Q)$ is finite. Put
\[
 U=\mathcal C(\mathcal C-1)H^0(Y_{\mathbb F_2},\Omega^1).
\]
Let $\mathcal S$ be a set of effective cuspidal divisors of degree $d$
over $\mathbb F_2$. If restriction from $U$ to every member of
$\mathcal S$ has rank $d$, no noncuspidal rational effective divisor
of degree $d$ on $Y$ has reduction in $\mathcal S$.
\end{lemma}
\begin{proof}
Suppose $D$ is such a divisor, and let $C$ be the unique etale
cuspidal lift of its reduction. The etale lift $C$ is fixed by the dyadic decomposition group.
It is therefore defined over a cyclotomic decomposition field whose
completion is $\mathbb Q_2$. Put $\alpha=T_2-a$ and $\beta=T_2-\delta$.
The class
\[
 \alpha[D-C]=\alpha[D-dc_0]-\alpha[C-dc_0]
\]
is torsion, by the finite-image hypothesis and Manin--Drinfeld,
and reduces to zero. Parent's lemma gives exponent two, and
Lemma~\ref{lem:nonic95-dyadic-Hecke} then gives
$\beta\alpha[D-C]=0$.

Consider $f:Y^{(d)}\longrightarrow J$ defined by
$E\longmapsto\beta\alpha[E-dc_0]$.
On the special fibre $T_2^*=\delta^*\mathcal C$, since Frobenius
pulls back differentials to zero. Since $a$ is even, the cotangent
image of $\beta\alpha$ is
\[
 (\delta^*)^2\mathcal C(\mathcal C-1)H^0(\Omega^1)=U.
\]
Here $\delta^*$ commutes with Cartier and is invertible.
The rank condition is precisely formal immersion at $\overline C$.
The sections $D,C$ have the same reduction and the same image under
$f$, hence coincide, a contradiction.
\end{proof}

\begin{remark}
If the only simple factors of $J$ with positive rank are isogenous
to an elliptic curve with even coefficient $a_2$, the finite-image
hypothesis holds with $a=a_2$. Thus the criterion can apply when
$J(\mathbb Q)$ has positive rank, without using a diamond difference.
On $J(\mathbb F_2)$, Eichler--Shimura also gives
$T_2=1+2\delta$. Consequently the reduction of $T_2-a$ on rational
divisor classes can be computed using $1-a+2\delta$.
\end{remark}

\begin{proposition}\label{prop:nonic-no77-continuation}
An elliptic curve over a number field of degree $9$ has no point of order $77$.
\end{proposition}
\begin{proof}
Put $X=X_1(77)$ and
$Y=X/\langle\langle36\rangle\rangle$, where the diamond subgroup
has order $5$. Its action is free, including at the cusps and in
characteristic $2$. Thus $Y$ has good reduction at $2$ and
$g(Y)=1+(181-1)/5=37$.

The fresh complete winding computation at level $77$ accounts for
all $181$ dimensions. Every nontrivial-character component is full
winding. The trivial-character component has dimension $7$, winding
 dimension $6$, and $T_2$ acts as zero on the one-dimensional quotient.
Hence Kato's theorem gives that $T_2J_1(77)(\mathbb Q)$ is finite.
Pullback and norm along $X\to Y$ show that $T_2J_Y(\mathbb Q)$ is
finite as well.

No elliptic curve over $\mathbb F_{2^e}$ has a point of order $77$
for $1\leq e\leq9$. Hasse's interval has no multiple of $77$ except
orders $77$, $231$, $539$ at $e=6,8,9$, respectively. Their traces
are $-12$, $26$, $-26$, all even. The supersingular classification
in \cite[Theorem~4.6]{DN} excludes them. This is checked independently
by the exact trace certificate.

Suppose a degree-nine torsion point exists, and let $D$ be its
rational effective trace divisor of degree $9$ on $X$.
Its reduction at $2$ is entirely cuspidal. The source cusp signature
is $1^{30},3^{10},10^3,30$.
Consequently its reduction contains only rational and cubic cusps.
Let $D_Y$ be its pushforward to $Y$. This is a noncuspidal rational
effective divisor of degree $9$.
The special-fibre cusp signature of $Y$, with widths, is
\[
 (77,1)^6,\quad(7,2)^3,\quad(11,3)^2,\quad(1,6).
\]
The rational source cusps map to the six rational cusps of width
$77$, and the cubic source cusps map to the two cubic cusps of
width $11$. Thus $\overline{D_Y}$ is supported on these eight cusps.
In particular, we retain the cubic cusp branch in degree $9$.

We recall the degree-independent modular identification of the
explicit model used in the computation. Over the elliptic curve
$X_0(11)_{\mathbb F_2}$ with its exact $j$-map, the irreducible equation
\[
 v^8+jv^7+v^4+1=0
\]
gives $X_0(77)_{\mathbb F_2}$ of genus $7$. The identity for
$j_7(v)=(v^2+v+1)^4/v^7$ and its dual is checked against $\Phi_7$.
The quotient $Y\to X_0(77)$ is an etale cyclic cover of degree $6$,
with cusp Frobenius values $0,2,3,5$ at widths $77,11,7,1$.
The base Picard group has invariants $30,60,0$.
Exactly six characters to $C_6$ have the required cusp values.
Two actual degree-ten places obtained from elliptic curves with
points of order $77$ must split in the cover and identify exactly
one of those characters.

Its cubic subcover is constructed as
$X_1(7)\times_{X(1)}X_0(11)$, using
\[
 b=t^3-t^2,\qquad c=t^2-t,\qquad
 v=\frac{t^2+t}{t^3+t+1}.
\]
The quadratic subcover is the fixed Artin--Schreier equation in
\texttt{AS\_component.m}. Its character is checked on places whose
Picard classes generate the base Picard group modulo $2$.
The composite is irreducible of genus $37$, has a rational cusp,
and has $j$-map degree $576$. Riemann--Hurwitz proves that the
cover is geometrically connected and unramified.
These modular-identification certificates are copied unchanged;
no earlier-degree exclusion is used. The model and both physical
deck actions are also rebuilt in the present degree-nine directory.

Apply Lemma~\ref{lem:nonic77-finite-image-Cartier} with $a=0$.
The relevant cotangent image is
\[
 U=\mathcal C(\mathcal C-1)H^0(Y_{\mathbb F_2},\Omega^1),
 \qquad \dim U=32.
\]
There are precisely
\[
 \binom{14}{5}+2\binom{11}{5}+3\binom85+4=3098
\]
effective divisors of degree $9$ on the six rational and two cubic
cusps. We verify that the restriction of $U$ has rank $9$ for all
of them.

The degree-independent rational-cusp Riemann--Roch kernels for
multiplicities up to $8$ are retained. To identify their coordinates
in the freshly built differential basis, we recompute Cartier,
the two deck matrices, and one multiplicity-eight kernel.
The linear intertwiner conditions have a one-dimensional solution
space, whose nonzero element is invertible. This identifies the
basis transition uniquely.
We then compute the multiplicity-nine kernel at one rational cusp
and the multiplicity-one, two and three kernels at one cubic cusp.
Verified deck transformations transport these kernels to the other
cusps. The resulting $60$ restriction blocks cover every allowed
multiplicity. Their kernels, the basis transition, the Cartier
projection and all $3098$ rank computations are replayed independently
over $\mathbb F_2$.

The restriction ranks give formal immersion at every possible
reduction of $D_Y$. The criterion therefore excludes $D_Y$, a
contradiction.
\end{proof}
